% ---------- study design through analysis plan ----------
\section{Study Design}

The experiment follows a factorial design with a longitudinal component at the check
level: 3 application tasks $\times$ 2 builders $\times$ 3 independent repetitions,
yielding 18 initial applications (iteration~0), each eligible for up to 3 repair
iterations. Figure~\ref{tab:design} summarizes the structure.

\begin{table}[tbp]
\centering
\caption{Design structure. Each cell contains one application per repetition; each
application is a sequence of at most four measured suite executions (i0--i3).}
\label{tab:design}
\small
\begin{tabularx}{\textwidth}{llX}
\toprule
Factor & Levels & Value \\
\midrule
Task & 3 & APP-01 booking; APP-02 dashboard; APP-03 inventory \\
Builder & 2 & two commercial builders meeting \S4 criteria \\
Repetition & 3 & independent generation runs per cell \\
Repair iterations & $\leq 3$ & whole-suite re-measurement after each \\
\bottomrule
\end{tabularx}
\end{table}

The scale is a deliberate trade-off. The statistical unit is the \emph{check
observation} (application $\times$ iteration $\times$ check); the frozen catalog
contains 39 checks, and per-application applicability and cascade rules yield
roughly 20--25 eligible checks per suite execution, hence roughly 1,500--2,000
eligible check observations across 72 suite executions and roughly 700--1,200
transition-eligible observations. At the
application level the sample is small, so all application-level
aggregates are reported with cluster-aware uncertainty, and the builder comparison
(RQ5) is treated as exploratory. Larger grids were rejected because commercial
builder subscriptions are metered, repair loops are operator-driven, and three
repetitions already permit within-cell variance estimation under generation
randomness.

\section{Benchmark and Application Tasks}

Three bounded applications were specified, each frozen in
\texttt{benchmark/specifications/}: an appointment booking system (two roles, CRUD,
search, admin panel), a multi-role SaaS dashboard (three roles, project--task
hierarchy, role-scoped API), and an ordering/inventory application (order placement
with stock decrement, manager-only inventory administration).

Each specification contains three components. First, numbered functional, security,
accessibility, and persistence requirements written in prose precise enough to be
checked. Second, a binding \emph{interface contract}: exact HTTP routes, request and
response JSON shapes, seed accounts, and DOM \texttt{data-testid} attributes. The
contract is what makes deterministic cross-implementation testing possible; builders
retain freedom over stack, architecture, and visual design, but generated
applications that deviate from the contract fail the affected checks, and that
failure is a measured outcome rather than a setup error. Third, an applicability map
to the frozen check catalog.

The prescriptive contract trades ecological representativeness for measurement
validity: freely prompted applications cannot be compared on deterministic oracles,
because the oracles would differ per application. We record this trade-off as a
threat in \S12 and note that it biases \emph{against} builders in the short term
(contract violations appear as failures) while favoring measurement consistency
across iterations.

\section{Evaluation Oracles}

The oracle suite contains 39 checks across eight categories
(Table~\ref{tab:oracles}), each emitted with a frozen identifier from the check
catalog. Suite execution is automated by a runner that installs, builds, launches,
and probes the application, and records a schema-validated run record with raw tool
output for every check.

\begin{table}[tbp]
\centering
\caption{Oracle categories. Runtime categories are recorded as not applicable when
the application cannot be started (cascade rule).}
\label{tab:oracles}
\setlength{\tabcolsep}{4.5pt}
\small
\begin{tabularx}{\textwidth}{llX}
\toprule
Category & Checks & Method \\
\midrule
BUILD & 5 & fresh \texttt{npm ci/install}; declared build; server start on \$PORT; root reachability; startup log hygiene \\
TYPE & 1 & \texttt{tsc --noEmit} for TypeScript projects \\
LINT & 2 & ESLint with a frozen framework-agnostic config; zero errors / zero warnings \\
FUNC & 10--11 & Playwright (Chromium) drives the interface contract: login, logout, create, list render, edit, delete, validation, search, empty state, not-found, API rejection \\
AUTHZ & 6 & API probes across the role $\times$ route matrix, cross-owner writes, UI denial, logout invalidation \\
SEC & 6 & frozen Semgrep ruleset; secret-pattern scan; \texttt{npm audit} high/critical; debug-endpoint probes; cookie flags or non-localStorage tokens; credential storage \\
A11Y & 4 & axe-core (WCAG~2.1 A/AA subset) on login, primary list, and form pages; language and title \\
DB & 4 & create, update, delete survive full process restart; response shapes match the contract \\
\bottomrule
\end{tabularx}
\end{table}

Three rules govern interpretation. \emph{Cascade}: when the application cannot be
started, runtime checks are recorded as not applicable rather than failed, and the
install/build failures carry the signal. \emph{Flakiness policy}: environment-class
failures (browser crash, localhost timeout) are retried once and logged; test-content
failures are never retried; during the pilot, every oracle was executed twice
back-to-back to detect non-deterministic checks. \emph{Security caution}: scanner
findings are candidate findings, not confirmed vulnerabilities. Every distinct
security finding is manually inspected before analysis; a sensitivity analysis
repeats the security results with findings judged false positives counted as passing.
No CVSS scores are assigned and no exploitable-vulnerability claims are made from
scanner output alone.

\section{Repair Protocol}

After each suite execution, a feedback generator renders one fixed-template message
per failing check---deterministic text with the observed and expected values and a
specification reference, e.g., \emph{``AUTHZ-003 failed: request GET /api/admin/stats
as role `patient' returned 200; expected 403.''}---preceded by a preamble instructing
the builder to repair without changing unrelated behavior. The operator delivers the
feedback verbatim to the same builder session that produced the application, exports
the resulting source, and the runner re-executes the \emph{entire} suite. The
whole-suite rule admits no exceptions: regressions are only observable to checks that
run, and rerunning only the failing subset would undercount them by construction.
Iteration~0 is archived once and never overwritten; repairs may stop early only when
no checks are failing.

\section{Analysis Plan}

\subsection{Metrics}

Per application $a$ and transition $k \to k{+}1$, with outcome function
$X_{a,k}(\cdot)$ over checks, let $T(a,k)=\{c: X_{a,k}(c)=\mathrm{fail}\}$ and
$P(a,k)=\{c: X_{a,k}(c)=\mathrm{pass}\}$, excluding environment-error and
not-applicable checks from all counts:
\begin{align}
\rsr(a,k) &= \frac{|\{c \in T(a,k) : X_{a,k+1}(c)=\mathrm{pass}\}|}{|T(a,k)|}, &
\rr(a,k) &= \frac{|\{c \in P(a,k) : X_{a,k+1}(c)=\mathrm{fail}\}|}{|P(a,k)|},\\
\nqc(a,k) &= |\text{repaired}| - |\text{regressed}|. && \nonumber
\end{align}
Raw counts accompany every rate; rates with fewer than five targets or five passing
checks are flagged as small-denominator. Because checks within an application are not
independent---a build failure cascades into every runtime check---rates are estimated
with application-level clustering, and a sensitivity analysis repeats category-level
comparisons after excluding cascading categories when BUILD fails at iteration $k$.

\subsection{Statistical procedures}

The plan is estimation-first. Descriptive trajectories per application and category
come first; uncertainty is quantified by a cluster bootstrap over applications
(2{,}000 replicates, percentile intervals), since checks are not independent units.
If a model is used, two transition datasets are modeled separately---fail$\to$pass
outcomes for RQ1 and pass$\to$fail outcomes for RQ2---with iteration $\times$
category fixed effects and random intercepts for application and check; a GLMM is
fitted for a transition type only if at least ten applications contribute at least
one event of that type, otherwise bootstrap-only inference is reported, and the
choice is stated. Two pre-declared hypothesis families---category contrasts (RQ3,
restricted to categories with at least eight eligible checks per application:
functionality, authorization, security, persistence) and iteration contrasts (RQ4)
---use Benjamini--Hochberg correction within family at $\alpha=.05$; smaller
categories (build, type, lint, accessibility) are reported descriptively. Effect
sizes are reported as risk differences with intervals. Engineering significance is
discussed separately from statistical significance, and no importance claim rests on
a p-value alone. All numbers in any future results section will be regenerated by
repository scripts from the processed dataset; none are transcribed by hand. Two
further points bound interpretation. Each repair transition bundles all failing
checks into a single builder action, so per-check rates estimate check-level
behavior under bundle feedback, not independent repair events. And the primary
estimands use transition-eligible sets (checks whose outcome is defined at both
iterations), with the protocol-literal fixed denominators reported as a conservative
sensitivity and an explicit exclusion ledger; the protocol changelog records this
amendment, made pre-data, alongside every other post-review change.

\section{Status of Results}

No experimental data has been collected at the time of this writing. The Results
section of the eventual manuscript does not exist yet, and this document does not
anticipate it. The protocol, specifications, prompts, oracles, and analysis plan are
frozen pre-data; every post-freeze change is recorded in the repository changelog,
and the changes below were all made before data collection, in response to
adversarial review passes.

What \emph{has} been executed is a harness-validation pilot, reported here only as
evidence that the machinery works and only under that label. The pilot generated the
APP-01 booking application with a non-study LLM system (GLM via a programmatic SDK;
it is not Lovable, Bolt, or Replit, and the run is excluded from all research data),
executed the full 39-check suite, delivered machine-generated feedback for sixteen
failures, executed a repair, re-ran the whole suite, and repeated once more. With
the corrected harness, the measured trajectory was 19 passing / 16 failing at
iteration~0; a repair that changed the code, repaired 0 of 16 failures, and
regressed 1 of 19 passing checks at iteration~1 (the builder moved the session token
into a cookie without the required security flags); and a correct repair of that
same cookie-flags check at iteration~2 (repair success 1/17, regression rate 0/18).
Persistent failures---missing interface attributes (which cascade into every
UI-driven check), no root route, a stateless-token logout that does not invalidate
sessions, a contract schema deviation, and one high and one critical dependency
vulnerability---are plausible model behavior and demonstrate precisely the kind of
dynamics the full study is designed to quantify. Fifteen harness defects found
during the pilot and review passes (among them a lint configuration that produced
spurious errors, an over-strict health check that misread redirects, orphaned server
processes, an invalid scanner configuration, silent pass-through of dependency-audit
failures, and cross-iteration state pollution) were fixed, and each fix is logged in
the protocol changelog and the revision log. The pilot's purpose was to find exactly
these defects before they could contaminate study data.
